% ============================================================================
%  展望  （原书 p.128，PDF p.131）
% ============================================================================
\clearpage
\section*{展望}
\addcontentsline{toc}{section}{展望}

在结束讨论之前，作一点回顾是应当的。同调镜像对称带我们走过了一段旅程：从大量"抽象废话"
（也就是范畴论与同调代数）出发，进入 Floer 理论所刻画的几何领域，最终又拱回
Kontsevich 的代数猜想，即我们最初声明的目标。这条抛物线要求我们向纯粹数学的许多基础
分支借用概念：辛几何与复几何、泛函分析的影子、无处不在的代数拓扑，最后还有代数几何的
一瞥——而这一切始终有范畴论的记号作伴。无需多言，我们才刚刚触及这个主题所能提供之物的
皮毛。而正是在这种看似永无止境的可能性中，也许就藏着它不可否认的魅力。

毕竟，一个尝试性的物理"万有理论"会汇聚如此丰富的数学、并指出其间的意外联系，从而提升
一种更微妙地交织在一起的理论图景——并反过来提升现实本身——这种愿景看来是公平的。当然，
据作者所能作出的最佳猜测，这里所呈现的部分内容已经相当过时；而同调镜像猜想至今在其最
一般的形式下仍未获证明。但教训仍在彼处，而那首交响乐值得一听。

% [Tutor Note: 【背景】"abstract nonsense" 是范畴论的戏称，源自 Mac Lane；作者用它自嘲，
% 中译保留为"抽象废话"并在括号内点明。末句 "the symphony is one worth listening to"
% 呼应前文 "symphony"，故译为"交响乐"，不译作"交响曲"以外的比喻义。]

% [Tutor Note: 【勘误】原书 p.128 正文首句 "a little retrospection is in order" 与末段
% "part of what presented here is ... already quite obsolete" 中，原文 "what presented"
% 缺系动词（应为 "what is presented"），属英文笔误，译文按语法补足，特此说明。]
